% =============================================================================
% Aula 08 — Thymeleaf e banco de dados {\textperiodcentered} Projeto banco
% Beamer + tema Madrid | Curso Entra21 — Spring Framework
% Compilar: tectonic -X compile slides-aula-08.tex
% =============================================================================
\documentclass[aspectratio=169,11pt]{beamer}

\usetheme{Madrid}
\usecolortheme{default}

\usepackage[T1]{fontenc}
\usepackage[utf8]{inputenc}
\usepackage[brazil]{babel}
\usepackage{booktabs}
\usepackage{graphicx}
\usepackage{hyperref}
\usepackage{listings}
\usepackage{xcolor}
\usepackage{multicol}
\usepackage{tikz}
\usetikzlibrary{arrows.meta, positioning, shapes.geometric}

\definecolor{codebg}{RGB}{245,245,245}
\definecolor{codegreen}{RGB}{40,120,40}
\definecolor{codeblue}{RGB}{30,70,160}
\definecolor{codered}{RGB}{180,50,50}

\lstset{
  basicstyle=\ttfamily\small,
  backgroundcolor=\color{codebg},
  frame=single,
  rulecolor=\color{black!20},
  breaklines=true,
  columns=fullflexible,
  keepspaces=true,
  showstringspaces=false,
  keywordstyle=\color{codeblue}\bfseries,
  commentstyle=\color{codegreen},
  stringstyle=\color{codered},
  tabsize=2,
  literate={á}{{\'a}}1 {é}{{\'e}}1 {í}{{\'i}}1 {ó}{{\'o}}1 {ú}{{\'u}}1
           {Á}{{\'A}}1 {É}{{\'E}}1 {Í}{{\'I}}1 {Ó}{{\'O}}1 {Ú}{{\'U}}1
           {ã}{{\~a}}1 {õ}{{\~o}}1 {Ã}{{\~A}}1 {Õ}{{\~O}}1
           {ç}{{\c{c}}}1 {Ç}{{\c{C}}}1
           {ê}{{\^e}}1 {Ê}{{\^E}}1 {â}{{\^a}}1 {Â}{{\^A}}1
           {ô}{{\^o}}1 {Ô}{{\^O}}1
           {à}{{\`a}}1 {À}{{\`A}}1
}

\lstdefinelanguage{props}{
  comment=[l]{\#},
  morestring=[b]"
}

\newcommand{\onde}[1]{%
  \begin{block}{Onde estamos}
    #1
  \end{block}
}

\title{Aula 08 --- Spring Framework}
\subtitle{Thymeleaf {\textperiodcentered} Banco de dados {\textperiodcentered} Projeto banco}
\author{Mauricio Duailibi Neto}
\institute{Turma Java Entra21}
\date{}

\begin{document}

\begin{frame}
  \titlepage
\end{frame}

\begin{frame}{Esta aula tem dois dias}
  \begin{columns}[T]
    \begin{column}{0.48\textwidth}
      \begin{block}{Dia 1 {\textperiodcentered} Parte 1 --- 18h30 às 20h}
        O que é Thymeleaf\\
        O que é banco de dados\\
        Instalar o XAMPP
      \end{block}
      \begin{block}{Dia 1 {\textperiodcentered} Parte 2 --- 20h20 às 22h}
        O projeto banco\\
        Configurar o banco no Spring\\
        Tela de cadastro gravando no MySQL
      \end{block}
    \end{column}
    \begin{column}{0.48\textwidth}
      \begin{block}{Dia 2 {\textperiodcentered} Parte 3 --- 18h30 às 20h}
        Login com CPF e senha\\
        Página da conta com o saldo\\
        Depósito e saque
      \end{block}
      \begin{block}{Dia 2 {\textperiodcentered} Parte 4 --- 20h20 às 22h}
        PIX para outra conta\\
        Extrato\\
        Teste com pessoas do gerador
      \end{block}
    \end{column}
  \end{columns}
\end{frame}

\begin{frame}{Duas novidades}
  Até a aula 07 a gente fazia sempre do mesmo jeito.\\
  A partir de hoje, duas coisas mudam.

  \vspace{0.6em}
  \begin{description}
    \item[Thymeleaf] O Java passa a montar a página HTML.\\
      Sem Axios. Sem \texttt{app.js}.
    \item[Banco de dados] Os dados deixam de ficar no \texttt{ArrayList}.\\
      Vão para o MySQL e não somem quando o Spring para.
  \end{description}

  \vspace{0.6em}
  O projeto que junta as duas coisas é um \textbf{banco}:\\
  conta, saldo, depósito, saque e PIX.
\end{frame}

\begin{frame}{Antes de começar: baixem o XAMPP}
  O download demora um pouco.\\
  Deixem baixando enquanto a gente conversa.

  \vspace{0.6em}
  \begin{enumerate}
    \item Abram \texttt{https://www.apachefriends.org}
    \item Cliquem em \textbf{XAMPP para Windows}
    \item Só baixem. \textbf{Não instalem ainda.}
  \end{enumerate}

  \vspace{0.6em}
  Daqui a pouco a gente instala junto, passo a passo.
\end{frame}

\begin{frame}{Mapa da aula}
  \begin{multicols}{2}
    \small
    \tableofcontents
  \end{multicols}
\end{frame}

% #############################################################################
\section{Como fazíamos até agora}

\begin{frame}{Relembrando o ATM da aula 07}
  No caixa eletrônico da aula passada:

  \vspace{0.5em}
  \begin{itemize}
    \item \texttt{index.html} ficava na pasta \texttt{static}
    \item \texttt{app.js} usava o Axios para chamar \texttt{/api/conta}
    \item O Java respondia com JSON: \texttt{\{"saldo": 50\}}
    \item O JavaScript pegava o número e colocava na tela
  \end{itemize}

  \vspace{0.5em}
  O Java só mandava \textbf{dados}.\\
  Quem montava a tela era o JavaScript, no navegador.
\end{frame}

\begin{frame}{O desenho de antes}
  \begin{center}
  \begin{tikzpicture}[
      box/.style={draw, rounded corners, align=center,
                  minimum width=2.8cm, minimum height=1.0cm,
                  font=\small},
      seta/.style={-{Latex}, thick}
    ]
    \node[box, fill=blue!15] (nav) {Navegador\\\scriptsize HTML + JS};
    \node[box, fill=green!20, right=4.5cm of nav] (java) {Spring\\\scriptsize \texttt{@RestController}};

    \draw[seta] ([yshift=0.25cm]nav.east) -- node[above, font=\scriptsize]
      {1. Axios: GET \texttt{/api/conta}} ([yshift=0.25cm]java.west);
    \draw[seta] ([yshift=-0.25cm]java.west) -- node[below, font=\scriptsize]
      {2. JSON: \texttt{\{"saldo": 50\}}} ([yshift=-0.25cm]nav.east);
  \end{tikzpicture}
  \end{center}

  \vspace{0.6em}
  3. O JavaScript pega o 50 e escreve \texttt{R\$ 50,00} na tela.

  \vspace{0.6em}
  Isso se chama \textbf{API REST}.\\
  Foi o que fizemos da aula 01 até a aula 07.
\end{frame}

\begin{frame}{O problema da memória}
  Onde ficavam os alunos, os produtos, as vendas?

  \vspace{0.5em}
  Num \texttt{ArrayList}, dentro da classe \texttt{Dados}.

  \vspace{0.5em}
  \begin{alertblock}{O que acontece quando o Spring para}
    Tudo some.\\
    Cadastrou dez produtos, parou o Spring, subiu de novo:\\
    a lista volta vazia.
  \end{alertblock}

  \vspace{0.5em}
  Num banco de verdade isso não pode acontecer.\\
  Ninguém aceita perder o saldo porque o servidor reiniciou.
\end{frame}

% #############################################################################
\section{O que é Thymeleaf}

\begin{frame}{Uma comparação com restaurante}
  \begin{block}{Antes (API REST)}
    A cozinha manda os ingredientes separados para a mesa:\\
    o arroz, o feijão, a carne.\\
    Quem monta o prato é o cliente (o JavaScript).
  \end{block}

  \vspace{0.5em}
  \begin{block}{Agora (Thymeleaf)}
    A cozinha já manda o \textbf{prato montado}.\\
    O cliente só come.
  \end{block}

  \vspace{0.5em}
  O ``prato montado'' é a página HTML com os dados já dentro.
\end{frame}

\begin{frame}{O que é o Thymeleaf}
  É uma ferramenta que o Spring usa para montar páginas HTML.

  \vspace{0.6em}
  A gente escreve um HTML normal,\\
  com alguns espaços marcados para o Java preencher.

  \vspace{0.6em}
  Antes de mandar a página para o navegador,\\
  o Thymeleaf troca esses espaços pelos dados de verdade.

  \vspace{0.6em}
  \begin{block}{Em uma frase}
    O servidor devolve a página pronta, já com o nome e o saldo.
  \end{block}
\end{frame}

\begin{frame}{O desenho de agora}
  \begin{center}
  \begin{tikzpicture}[
      box/.style={draw, rounded corners, align=center,
                  minimum width=2.6cm, minimum height=1.0cm,
                  font=\small},
      seta/.style={-{Latex}, thick}
    ]
    \node[box, fill=blue!15] (nav) {Navegador\\\scriptsize só HTML};
    \node[box, fill=green!20, right=3.6cm of nav] (java) {Spring\\\scriptsize \texttt{@Controller}};
    \node[box, fill=yellow!25, right=1.6cm of java] (tpl) {Thymeleaf\\\scriptsize \texttt{conta.html}};

    \draw[seta] ([yshift=0.25cm]nav.east) -- node[above, font=\scriptsize]
      {1. GET \texttt{/conta}} ([yshift=0.25cm]java.west);
    \draw[seta] ([yshift=0.25cm]java.east) -- node[above, font=\scriptsize]
      {2. dados} ([yshift=0.25cm]tpl.west);
    \draw[seta] ([yshift=-0.25cm]tpl.west) -- node[below, font=\scriptsize]
      {3. HTML} ([yshift=-0.25cm]java.east);
    \draw[seta] ([yshift=-0.25cm]java.west) -- node[below, font=\scriptsize]
      {4. página pronta} ([yshift=-0.25cm]nav.east);
  \end{tikzpicture}
  \end{center}

  \vspace{0.6em}
  O navegador recebe \texttt{<p>R\$ 50,00</p>} já escrito.\\
  Não tem Axios. Não tem JavaScript para montar a tela.
\end{frame}

\begin{frame}[fragile]{Um exemplo pequeno}
  No arquivo HTML a gente escreve:

\begin{lstlisting}[language=HTML,basicstyle=\ttfamily\footnotesize]
<p th:text="${nome}">aqui vai o nome</p>
\end{lstlisting}

  No Java a gente diz qual é o nome:

\begin{lstlisting}[language=Java,basicstyle=\ttfamily\footnotesize]
model.addAttribute("nome", "Ana");
\end{lstlisting}

  O navegador recebe:

\begin{lstlisting}[language=HTML,basicstyle=\ttfamily\footnotesize]
<p>Ana</p>
\end{lstlisting}

  O \texttt{th:text} trocou o texto de dentro pelo valor do Java.
\end{frame}

\begin{frame}{Antes e agora, lado a lado}
  \begin{center}
  \begin{tabular}{lll}
    \toprule
     & Antes (REST) & Agora (Thymeleaf) \\
    \midrule
    Anotação & \texttt{@RestController} & \texttt{@Controller} \\
    O Java devolve & dados (JSON) & o nome da página \\
    Quem monta a tela & JavaScript & o servidor \\
    Como o form envia & Axios & o próprio \texttt{<form>} \\
    Pasta do HTML & \texttt{static/} & \texttt{templates/} \\
    Arquivo \texttt{app.js} & sim & não precisa \\
    \bottomrule
  \end{tabular}
  \end{center}

  \vspace{0.6em}
  A pasta \texttt{static/} continua existindo.\\
  Ela guarda o CSS e as imagens.
\end{frame}

\begin{frame}[fragile]{A diferença no Controller}
  \begin{columns}[T]
    \begin{column}{0.48\textwidth}
      \textbf{Antes}
\begin{lstlisting}[language=Java,basicstyle=\ttfamily\scriptsize]
@RestController
public class ContaController {

  @GetMapping("/api/conta")
  public Conta ver() {
    return conta;
  }
}
\end{lstlisting}
      Devolve o objeto.\\
      Vira JSON.
    \end{column}
    \begin{column}{0.48\textwidth}
      \textbf{Agora}
\begin{lstlisting}[language=Java,basicstyle=\ttfamily\scriptsize]
@Controller
public class ContaController {

  @GetMapping("/conta")
  public String ver(Model model) {
    model.addAttribute("conta", conta);
    return "conta";
  }
}
\end{lstlisting}
      Devolve o texto \texttt{"conta"}.\\
      O Spring abre \texttt{templates/conta.html}.
    \end{column}
  \end{columns}
\end{frame}

\begin{frame}{Qual dos dois é melhor?}
  Nenhum. São usados em situações diferentes.

  \vspace{0.6em}
  \begin{description}
    \item[REST] Quando a tela é separada do servidor:\\
      aplicativo de celular, site feito em React, outro sistema.
    \item[Thymeleaf] Quando o site inteiro fica num projeto só.\\
      Muito sistema interno de empresa é assim.
  \end{description}

  \vspace{0.6em}
  Os dois são Spring.\\
  Quem sabe os dois consegue trabalhar em qualquer um.
\end{frame}

% #############################################################################
\section{A sintaxe do Thymeleaf}

\begin{frame}[fragile]{A primeira linha muda}
  Todo arquivo que usa Thymeleaf começa assim:

\begin{lstlisting}[language=HTML,basicstyle=\ttfamily\footnotesize]
<!DOCTYPE html>
<html lang="pt-BR" xmlns:th="http://www.thymeleaf.org">
\end{lstlisting}

  \vspace{0.4em}
  O \texttt{xmlns:th} avisa: ``neste arquivo existem comandos \texttt{th:}''.

  \vspace{0.5em}
  Sem essa linha a página até funciona,\\
  mas o VS Code sublinha tudo que começa com \texttt{th:}.
\end{frame}

\begin{frame}[fragile]{O Model é uma bandeja}
  O Controller coloca os dados numa bandeja.\\
  A página pega os dados da bandeja.

\begin{lstlisting}[language=Java,basicstyle=\ttfamily\footnotesize]
@GetMapping("/conta")
public String ver(Model model) {
  model.addAttribute("nome", "Ana");
  model.addAttribute("saldo", 250.0);
  return "conta";
}
\end{lstlisting}

  \texttt{"nome"} é a etiqueta.\\
  \texttt{"Ana"} é o que está dentro.\\
  Na página, a gente chama pela etiqueta: \texttt{\$\{nome\}}.
\end{frame}

\begin{frame}[fragile]{th:text --- escrever um valor}
\begin{lstlisting}[language=HTML,basicstyle=\ttfamily\footnotesize]
<p th:text="${nome}">Nome</p>
<p th:text="'Saldo: R$ ' + ${saldo}">R$ 0</p>
\end{lstlisting}

  \vspace{0.4em}
  \begin{itemize}
    \item \texttt{\$\{...\}} busca no Model
    \item O texto entre as tags (\texttt{Nome}) é trocado
    \item Texto fixo vai entre aspas simples: \texttt{'Saldo: R\$ '}
    \item O \texttt{+} junta os pedaços, igual no Java
  \end{itemize}

  \vspace{0.4em}
  Se o objeto tiver atributos, a gente usa o ponto:\\
  \texttt{\$\{conta.nome\}} chama o \texttt{getNome()} da conta.
\end{frame}

\begin{frame}[fragile]{th:if --- mostrar só às vezes}
\begin{lstlisting}[language=HTML,basicstyle=\ttfamily\footnotesize]
<p class="alerta-erro" th:if="${erro}" th:text="${erro}"></p>
\end{lstlisting}

  \vspace{0.4em}
  Se existir um \texttt{erro} no Model, o parágrafo aparece.\\
  Se não existir, o parágrafo nem vai para o navegador.

  \vspace{0.5em}
  É o \texttt{if} do Java, dentro do HTML.

  \vspace{0.5em}
  Vamos usar para as mensagens:\\
  ``Saldo insuficiente'', ``CPF ou senha incorretos''.
\end{frame}

\begin{frame}[fragile]{th:each --- repetir para cada item}
\begin{lstlisting}[language=HTML,basicstyle=\ttfamily\footnotesize]
<ul>
  <li th:each="nome : ${nomes}" th:text="${nome}"></li>
</ul>
\end{lstlisting}

  Se \texttt{nomes} for uma lista com Ana, Bruno e Carla, sai:

\begin{lstlisting}[language=HTML,basicstyle=\ttfamily\footnotesize]
<ul>
  <li>Ana</li>
  <li>Bruno</li>
  <li>Carla</li>
</ul>
\end{lstlisting}

  É o \texttt{for} do Java.\\
  Vamos usar no extrato: uma linha para cada movimentação.
\end{frame}

\begin{frame}[fragile]{th:action e th:href --- os endereços}
\begin{lstlisting}[language=HTML,basicstyle=\ttfamily\footnotesize]
<form th:action="@{/deposito}" method="post">
<a th:href="@{/sair}">Sair</a>
\end{lstlisting}

  \vspace{0.4em}
  \texttt{\$\{...\}} é para \textbf{dados}.\\
  \texttt{@\{...\}} é para \textbf{endereços}.

  \vspace{0.5em}
  O \texttt{th:action} diz para onde o formulário vai.\\
  O \texttt{th:href} diz para onde o link vai.

  \vspace{0.5em}
  Funcionaria com \texttt{action} e \texttt{href} comuns também.\\
  Usamos a versão \texttt{th:} porque é o jeito do Thymeleaf.
\end{frame}

\begin{frame}[fragile]{O formulário volta a trabalhar sozinho}
  Sem Axios, quem envia é o próprio \texttt{<form>}.

\begin{lstlisting}[language=HTML,basicstyle=\ttfamily\footnotesize]
<form th:action="@{/deposito}" method="post">
  <input type="number" name="valor">
  <button type="submit">Depositar</button>
</form>
\end{lstlisting}

  No Java, o \texttt{name} do input vira o parâmetro:

\begin{lstlisting}[language=Java,basicstyle=\ttfamily\footnotesize]
@PostMapping("/deposito")
public String depositar(@RequestParam double valor) {
\end{lstlisting}

  \texttt{name="valor"} no HTML, \texttt{valor} no Java.\\
  Se os nomes forem diferentes, dá erro.
\end{frame}

\begin{frame}[fragile]{redirect --- mandar o navegador para outro lugar}
\begin{lstlisting}[language=Java,basicstyle=\ttfamily\footnotesize]
return "conta";            // mostra templates/conta.html
return "redirect:/conta";  // manda o navegador abrir /conta
\end{lstlisting}

  \vspace{0.4em}
  Depois de um POST, a gente sempre usa \texttt{redirect}.

  \vspace{0.4em}
  \begin{block}{Por quê?}
    Sem o redirect, se a pessoa apertar F5 depois de depositar,\\
    o navegador manda o depósito de novo.\\
    Com o redirect, o F5 só recarrega a página da conta.
  \end{block}
\end{frame}

\begin{frame}{O navegador nunca vê o th:}
  Cliquem com o botão direito na página e em \textbf{Exibir código-fonte}.

  \vspace{0.6em}
  Não aparece \texttt{th:text}, nem \texttt{\$\{conta.nome\}}.\\
  Aparece o HTML pronto, com o nome escrito.

  \vspace{0.6em}
  Todo o trabalho do Thymeleaf acontece no servidor,\\
  antes da página sair.
\end{frame}

\begin{frame}{Resumo da sintaxe}
  \begin{center}
  \begin{tabular}{lll}
    \toprule
    Comando & Para que serve & Parecido com \\
    \midrule
    \texttt{th:text} & escrever um valor & \texttt{System.out.println} \\
    \texttt{th:if} & mostrar só se for verdade & \texttt{if} \\
    \texttt{th:each} & repetir para cada item & \texttt{for} \\
    \texttt{th:action} & para onde o form vai & URL do Axios \\
    \texttt{th:href} & para onde o link vai & \texttt{href} \\
    \texttt{\$\{...\}} & pegar dado do Model & variável \\
    \texttt{@\{...\}} & montar um endereço & \texttt{"/api/..."} \\
    \bottomrule
  \end{tabular}
  \end{center}

  \vspace{0.5em}
  É só isso que vamos usar nessa aula.
\end{frame}

% #############################################################################
\section{Banco de dados}

\begin{frame}{O que é um banco de dados}
  Um programa que guarda informação no disco, de forma organizada.

  \vspace{0.5em}
  Pensem numa planilha do Excel:

  \vspace{0.4em}
  \begin{center}
  \begin{tabular}{ll}
    \toprule
    Excel & Banco de dados \\
    \midrule
    O arquivo & o \textbf{banco} \\
    Cada aba & uma \textbf{tabela} \\
    Cada coluna & uma \textbf{coluna} (nome, cpf, saldo) \\
    Cada linha & um \textbf{registro} (uma pessoa) \\
    \bottomrule
  \end{tabular}
  \end{center}

  \vspace{0.5em}
  Desligou o computador? Os dados continuam lá.
\end{frame}

\begin{frame}{Uma tabela de verdade}
  A tabela \texttt{conta}, que vamos criar hoje:

  \vspace{0.5em}
  \begin{center}
  \begin{tabular}{lllll}
    \toprule
    id & nome & cpf & senha & saldo \\
    \midrule
    1 & Ana Souza & 11111111111 & 123 & 250.25 \\
    2 & Bruno Lima & 22222222222 & 456 & 150.25 \\
    3 & Carla Dias & 33333333333 & 789 & 0.00 \\
    \bottomrule
  \end{tabular}
  \end{center}

  \vspace{0.5em}
  Cada linha é uma conta.\\
  É igual ao \texttt{ArrayList<Conta>}, só que guardado no disco.
\end{frame}

\begin{frame}{A chave primária}
  A coluna \texttt{id} é especial.

  \vspace{0.4em}
  \begin{itemize}
    \item Nunca se repete
    \item Nunca fica vazia
    \item O próprio banco escolhe o número: 1, 2, 3\ldots
  \end{itemize}

  \vspace{0.5em}
  Ela se chama \textbf{chave primária}.

  \vspace{0.5em}
  Lembram do \texttt{Dados.proximoId = Dados.proximoId + 1}?\\
  O banco faz isso sozinho.
\end{frame}

\begin{frame}{MySQL e SQL}
  \begin{description}
    \item[MySQL] Um dos bancos de dados mais usados no mundo.\\
      É gratuito.
    \item[SQL] A língua que se usa para falar com o banco.\\
      Serve para o MySQL e para quase todos os outros.
  \end{description}

  \vspace{0.6em}
  O SQL tem quatro comandos principais.\\
  Vocês já conhecem a ideia.
\end{frame}

\begin{frame}{Os quatro comandos}
  \begin{center}
  \begin{tabular}{lll}
    \toprule
    SQL & O que faz & Parecido com \\
    \midrule
    \texttt{SELECT} & ler & GET \\
    \texttt{INSERT} & criar & POST \\
    \texttt{UPDATE} & alterar & PUT \\
    \texttt{DELETE} & apagar & DELETE \\
    \bottomrule
  \end{tabular}
  \end{center}

  \vspace{0.6em}
  É o mesmo CRUD das aulas 03 e 04.\\
  Só que agora falando direto com o banco.
\end{frame}

\begin{frame}[fragile]{Como o SQL se parece}
\begin{lstlisting}[language=SQL,basicstyle=\ttfamily\footnotesize]
SELECT * FROM conta;

SELECT nome, saldo FROM conta WHERE cpf = '11111111111';

INSERT INTO conta (nome, cpf, senha, saldo)
VALUES ('Ana Souza', '11111111111', '123', 0);

UPDATE conta SET saldo = 500 WHERE id = 1;
\end{lstlisting}

  \vspace{0.3em}
  Dá para ler quase como português:\\
  ``selecione nome e saldo da conta onde o cpf é\ldots''
\end{frame}

\begin{frame}{Vamos precisar escrever SQL?}
  Quase nada.

  \vspace{0.6em}
  O Spring escreve o SQL por nós.\\
  A gente chama um método Java e ele monta o comando.

  \vspace{0.6em}
  Vamos usar o SQL só para \textbf{olhar} o banco:\\
  conferir se a conta foi gravada, se o saldo mudou.

  \vspace{0.6em}
  Para isso existe o phpMyAdmin.
\end{frame}

% #############################################################################
\section{XAMPP e phpMyAdmin}

\begin{frame}{O que é o XAMPP}
  Um instalador que traz várias ferramentas juntas.

  \vspace{0.5em}
  Das que vêm nele, vamos usar três:

  \vspace{0.4em}
  \begin{description}
    \item[MySQL] O banco de dados.\\
      {\small (O XAMPP traz o MariaDB, que é o MySQL com outro nome.)}
    \item[phpMyAdmin] Uma página para ver e mexer no banco pelo navegador.
    \item[Apache] Só serve para abrir o phpMyAdmin.
  \end{description}

  \vspace{0.5em}
  O nosso Spring \textbf{não} usa o Apache.\\
  O Spring conversa direto com o MySQL.
\end{frame}

\begin{frame}{Passo 1 --- Instalar o XAMPP}
  \begin{enumerate}
    \item Abram o arquivo que baixaram
    \item Se o Windows perguntar se permite alterações: \textbf{Sim}
    \item Se aparecer um aviso sobre antivírus ou UAC: \textbf{OK}
    \item Nos componentes, deixem como está e cliquem em \textbf{Next}
    \item Pasta: deixem \texttt{C:\textbackslash xampp}
    \item \textbf{Next} até terminar
  \end{enumerate}

  \vspace{0.5em}
  \begin{alertblock}{Firewall}
    Se o Windows perguntar sobre o firewall, cliquem em \textbf{Permitir}.
  \end{alertblock}
\end{frame}

\begin{frame}{Passo 2 --- Ligar o MySQL}
  Abram o \textbf{XAMPP Control Panel}.

  \vspace{0.5em}
  \begin{enumerate}
    \item Na linha \textbf{Apache}, cliquem em \textbf{Start}
    \item Na linha \textbf{MySQL}, cliquem em \textbf{Start}
    \item Os dois nomes ficam com fundo verde
    \item Na linha do MySQL aparece a porta \texttt{3306}
  \end{enumerate}

  \vspace{0.5em}
  \begin{block}{Todo dia de aula}
    O MySQL precisa estar ligado \textbf{antes} de subir o Spring.\\
    Sem ele, o Spring não consegue iniciar.
  \end{block}
\end{frame}

\begin{frame}{Passo 3 --- Abrir o phpMyAdmin}
  No navegador:

  \vspace{0.4em}
  \begin{center}
    \large\texttt{http://localhost/phpmyadmin}
  \end{center}

  \vspace{0.5em}
  Do lado esquerdo aparece a lista de bancos que já existem.\\
  Alguns vêm com o XAMPP: \texttt{mysql}, \texttt{phpmyadmin}, \texttt{test}.

  \vspace{0.5em}
  Não mexam neles.\\
  Vamos criar o nosso.
\end{frame}

\begin{frame}{Passo 4 --- Criar o banco}
  \begin{enumerate}
    \item No lado esquerdo, cliquem em \textbf{Novo} (ou \textbf{New})
    \item Nome do banco: \texttt{banco}
    \item Ao lado, escolham \texttt{utf8mb4\_general\_ci}
    \item Cliquem em \textbf{Criar} (ou \textbf{Create})
  \end{enumerate}

  \vspace{0.5em}
  O banco \texttt{banco} aparece na lista da esquerda.\\
  Ainda sem nenhuma tabela.

  \vspace{0.5em}
  O \texttt{utf8mb4} deixa gravar acento e cedilha sem virar símbolo estranho.
\end{frame}

\begin{frame}[fragile]{Outro jeito: pela aba SQL}
  Quem preferir, pode criar escrevendo o comando.

  \vspace{0.4em}
  Cliquem na aba \textbf{SQL}, colem e cliquem em \textbf{Executar}:

\begin{lstlisting}[language=SQL,basicstyle=\ttfamily\footnotesize]
CREATE DATABASE banco
  CHARACTER SET utf8mb4
  COLLATE utf8mb4_general_ci;
\end{lstlisting}

  \vspace{0.4em}
  O resultado é o mesmo do passo 4.\\
  Não precisa fazer os dois.
\end{frame}

\begin{frame}{Usuário e senha do banco}
  O MySQL do XAMPP vem com:

  \vspace{0.4em}
  \begin{itemize}
    \item usuário: \texttt{root}
    \item senha: nenhuma, em branco
  \end{itemize}

  \vspace{0.5em}
  Isso só serve no nosso computador, para estudar.\\
  Um banco de verdade, na internet, sempre tem senha forte.

  \vspace{0.5em}
  Vamos colocar esse usuário na configuração do Spring daqui a pouco.
\end{frame}

\begin{frame}{Se o XAMPP não ligar}
  \begin{itemize}
    \item MySQL não fica verde --- outro MySQL já está instalado
          no computador usando a porta 3306. Chamem o professor.
    \item Apache não fica verde --- a porta 80 está ocupada
          (Skype, IIS). O phpMyAdmin não abre, mas o Spring funciona.
    \item \texttt{MySQL shutdown unexpectedly} --- chamem o professor.
    \item phpMyAdmin não abre --- confiram se o Apache está verde.
  \end{itemize}
\end{frame}

\begin{frame}{Checkpoint}
  \begin{itemize}
    \item [ ] XAMPP instalado
    \item [ ] Apache e MySQL verdes no painel
    \item [ ] \texttt{http://localhost/phpmyadmin} abre
    \item [ ] O banco \texttt{banco} aparece na lista da esquerda
  \end{itemize}

  \vspace{0.6em}
  Todo mundo com os quatro itens?\\
  Então vamos ver como o Java fala com esse banco.
\end{frame}

% #############################################################################
\section{JPA: o Java fala com o banco}

\begin{frame}{Duas línguas diferentes}
  O Java pensa em \textbf{objetos}: uma \texttt{Conta} com nome e saldo.

  \vspace{0.4em}
  O banco pensa em \textbf{tabelas}: linhas e colunas.

  \vspace{0.6em}
  Alguém precisa traduzir.\\
  Esse tradutor se chama \textbf{JPA}.

  \vspace{0.6em}
  \begin{center}
  \begin{tabular}{ll}
    \toprule
    Java & Banco \\
    \midrule
    classe \texttt{Conta} & tabela \texttt{conta} \\
    atributo \texttt{saldo} & coluna \texttt{saldo} \\
    um objeto \texttt{Conta} & uma linha da tabela \\
    \bottomrule
  \end{tabular}
  \end{center}
\end{frame}

\begin{frame}{Quem conversa com quem}
  \begin{center}
  \begin{tikzpicture}[
      box/.style={draw, rounded corners, align=center,
                  minimum width=2.4cm, minimum height=1.0cm,
                  font=\small},
      seta/.style={-{Latex}, thick}
    ]
    \node[box, fill=blue!15] (nav) {Navegador};
    \node[box, fill=green!20, right=1.2cm of nav] (ctl) {Controller};
    \node[box, fill=yellow!25, right=1.2cm of ctl] (rep) {Repository};
    \node[box, fill=orange!25, right=1.2cm of rep] (db) {MySQL};

    \draw[seta] (nav) -- (ctl);
    \draw[seta] (ctl) -- (rep);
    \draw[seta] (rep) -- (db);
  \end{tikzpicture}
  \end{center}

  \vspace{0.6em}
  \begin{itemize}
    \item O \textbf{Controller} recebe o pedido da tela
    \item O \textbf{Repository} é quem fala com o banco
    \item O Controller nunca escreve SQL. Ele pede para o Repository.
  \end{itemize}
\end{frame}

\begin{frame}[fragile]{Anotações na classe}
\begin{lstlisting}[language=Java,basicstyle=\ttfamily\footnotesize]
@Entity
public class Conta {

  @Id
  @GeneratedValue(strategy = GenerationType.IDENTITY)
  private Long id;

  @Column(unique = true)
  private String cpf;
\end{lstlisting}

  \begin{itemize}
    \item \texttt{@Entity} --- esta classe vira uma tabela
    \item \texttt{@Id} --- este atributo é a chave primária
    \item \texttt{@GeneratedValue} --- o banco escolhe o número
    \item \texttt{@Column(unique = true)} --- não pode repetir
  \end{itemize}
\end{frame}

\begin{frame}{Por que Long e não int?}
  Até agora o id era \texttt{int}.\\
  Com banco de dados a gente usa \texttt{Long}.

  \vspace{0.5em}
  \begin{itemize}
    \item \texttt{Long} é um número inteiro, que cabe valores bem grandes
    \item Com L maiúsculo ele pode ficar \textbf{vazio} (\texttt{null})
    \item Antes de salvar, a conta ainda não tem id: fica \texttt{null}
    \item Depois de salvar, o banco preenche: 1, 2, 3\ldots
  \end{itemize}

  \vspace{0.5em}
  É o costume em todo projeto Spring com banco.
\end{frame}

\begin{frame}[fragile]{O Repository}
\begin{lstlisting}[language=Java,basicstyle=\ttfamily\footnotesize]
public interface ContaRepository
    extends JpaRepository<Conta, Long> {
}
\end{lstlisting}

  \vspace{0.4em}
  É uma \textbf{interface}, não uma classe.\\
  A gente não escreve o código dos métodos.

  \vspace{0.4em}
  O Spring lê \texttt{<Conta, Long>} e entende:\\
  ``é um repositório de \texttt{Conta}, e o id é \texttt{Long}''.

  \vspace{0.4em}
  Ao subir, o Spring cria sozinho a classe que faz o trabalho.
\end{frame}

\begin{frame}{Métodos que já vêm prontos}
  \begin{center}
  \begin{tabular}{ll}
    \toprule
    Antes, com \texttt{Dados} & Agora, com o Repository \\
    \midrule
    \texttt{Dados.contas.add(conta)} & \texttt{contaRepository.save(conta)} \\
    \texttt{Dados.contas} & \texttt{contaRepository.findAll()} \\
    \texttt{for} procurando o id & \texttt{contaRepository.findById(id)} \\
    \texttt{Dados.contas.remove(...)} & \texttt{contaRepository.deleteById(id)} \\
    \bottomrule
  \end{tabular}
  \end{center}

  \vspace{0.6em}
  O \texttt{save} serve para os dois casos:\\
  se a conta é nova, faz \texttt{INSERT}.\\
  Se já existe, faz \texttt{UPDATE}.
\end{frame}

\begin{frame}[fragile]{Um método que a gente inventa}
  Para o login, precisamos achar a conta pelo CPF.

\begin{lstlisting}[language=Java,basicstyle=\ttfamily\footnotesize]
Conta findByCpf(String cpf);
\end{lstlisting}

  \vspace{0.3em}
  Só escrevemos essa linha, sem corpo.\\
  O Spring lê o \textbf{nome} do método:

  \vspace{0.3em}
  \begin{center}
  \begin{tabular}{ll}
    \toprule
    Pedaço & Significa \\
    \midrule
    \texttt{find} & buscar \\
    \texttt{By} & onde \\
    \texttt{Cpf} & o atributo \texttt{cpf} é igual ao que eu passar \\
    \bottomrule
  \end{tabular}
  \end{center}

  \vspace{0.3em}
  E monta: \texttt{SELECT * FROM conta WHERE cpf = ?}\\
  Se não achar, devolve \texttt{null}.
\end{frame}

\begin{frame}[fragile]{@Autowired --- pedir um pronto}
\begin{lstlisting}[language=Java,basicstyle=\ttfamily\footnotesize]
@Controller
public class AcessoController {

  @Autowired
  private ContaRepository contaRepository;
\end{lstlisting}

  \vspace{0.4em}
  A gente não faz \texttt{new ContaRepository()}.\\
  Nem daria: é uma interface.

  \vspace{0.4em}
  O \texttt{@Autowired} diz ao Spring:\\
  ``me entrega aquele repositório que você criou''.
\end{frame}

\begin{frame}{E a tabela, quem cria?}
  Também o Spring.

  \vspace{0.5em}
  Na configuração vamos colocar:\\
  \texttt{spring.jpa.hibernate.ddl-auto=update}

  \vspace{0.5em}
  Quando o Spring sobe, ele olha as classes com \texttt{@Entity}:
  \begin{itemize}
    \item Se a tabela não existe, ele cria
    \item Se faltar uma coluna, ele acrescenta
    \item Os dados que já estão lá, ele não apaga
  \end{itemize}

  \vspace{0.5em}
  A gente só cria o \textbf{banco} no phpMyAdmin.\\
  As \textbf{tabelas} ficam por conta do Spring.
\end{frame}

\begin{frame}{Resumo da teoria}
  \begin{enumerate}
    \item Thymeleaf: o servidor devolve a página pronta
    \item \texttt{@Controller} devolve o nome do HTML, que fica em \texttt{templates/}
    \item O Model leva os dados para a página
    \item Banco de dados guarda no disco. Nada some ao reiniciar.
    \item \texttt{@Entity} vira tabela, o Repository fala com o banco
    \item O XAMPP liga o MySQL. O phpMyAdmin mostra o que tem dentro.
  \end{enumerate}
\end{frame}

\begin{frame}{Intervalo}
  \begin{center}
    \Large Retomamos às \textbf{20h20}\\[0.8em]
    \normalsize
    Deixem o XAMPP aberto, com o MySQL ligado.\\
    Na volta a gente abre o projeto \textbf{banco}.
  \end{center}
\end{frame}

% #############################################################################
\section{O projeto banco}

\begin{frame}{O que vamos construir}
  O \textbf{Banco Entra21}. Nele, a pessoa pode:

  \vspace{0.5em}
  \begin{enumerate}
    \item Abrir uma conta com nome, CPF e senha
    \item Entrar com CPF e senha
    \item Ver o próprio saldo
    \item Depositar
    \item Sacar, se tiver saldo
    \item Mandar um PIX para outra conta, usando o CPF
    \item Ver o extrato de tudo que aconteceu
  \end{enumerate}

  \vspace{0.5em}
  É o ATM da aula 07, só que com várias contas\\
  e com o dinheiro guardado no banco de dados.
\end{frame}

\begin{frame}{As três telas}
  \begin{center}
  \begin{tikzpicture}[
      tela/.style={draw, rounded corners, align=left, fill=white,
                   minimum width=3.6cm, minimum height=3.0cm,
                   font=\scriptsize, anchor=north},
      titulo/.style={font=\small\bfseries},
      seta/.style={-{Latex}, thick}
    ]
    \node[titulo] (t1) at (0,0) {cadastro.html};
    \node[tela, below=0.1cm of t1] (c) {%
      Abrir conta\\[0.3em]
      Nome \hfill \fbox{\phantom{xxxxxx}}\\
      CPF \hfill \fbox{\phantom{xxxxxx}}\\
      Senha \hfill \fbox{\phantom{xxxxxx}}\\[0.3em]
      \colorbox{blue!20}{Cadastrar}\\[0.2em]
      \underline{Já tenho conta}};

    \node[titulo] (t2) at (4.8,0) {login.html};
    \node[tela, below=0.1cm of t2] (l) {%
      Entrar\\[0.3em]
      CPF \hfill \fbox{\phantom{xxxxxx}}\\
      Senha \hfill \fbox{\phantom{xxxxxx}}\\[0.3em]
      \colorbox{blue!20}{Entrar}\\[0.2em]
      \underline{Abrir uma conta}};

    \node[titulo] (t3) at (9.6,0) {conta.html};
    \node[tela, below=0.1cm of t3] (k) {%
      Olá, Ana \hfill \underline{Sair}\\[0.2em]
      \colorbox{blue!20}{Saldo R\$ 250,25}\\[0.2em]
      Depósito {\textperiodcentered} Saque {\textperiodcentered} PIX\\[0.2em]
      Extrato\\
      \quad Depósito \hfill 500,50\\
      \quad Saque \hfill -100,00};

    \draw[seta] (c.east) -- (l.west);
    \draw[seta] (l.east) -- (k.west);
  \end{tikzpicture}
  \end{center}
\end{frame}

\begin{frame}{O caminho de quem usa}
  \begin{enumerate}
    \item Abre \texttt{/cadastro}, preenche e clica em Cadastrar
    \item O Java grava a conta no banco, com saldo zero
    \item Vai para \texttt{/login}, digita CPF e senha
    \item O Java confere no banco e guarda quem entrou
    \item Vai para \texttt{/conta}: nome, saldo e as operações
    \item Cada operação é um formulário que volta para \texttt{/conta}
    \item \texttt{/sair} esquece quem estava logado
  \end{enumerate}
\end{frame}

\begin{frame}{Todos os endereços}
  \begin{center}
  \small
  \begin{tabular}{lll}
    \toprule
    Verbo e endereço & O que manda & O que acontece \\
    \midrule
    GET \texttt{/cadastro} & --- & mostra a tela de cadastro \\
    POST \texttt{/cadastro} & nome, cpf, senha & grava a conta \\
    GET \texttt{/login} & --- & mostra a tela de login \\
    POST \texttt{/login} & cpf, senha & confere e entra \\
    GET \texttt{/conta} & --- & mostra saldo e extrato \\
    POST \texttt{/deposito} & valor & soma no saldo \\
    POST \texttt{/saque} & valor & tira do saldo \\
    POST \texttt{/pix} & cpfDestino, valor & passa para outra conta \\
    GET \texttt{/sair} & --- & sai da conta \\
    \bottomrule
  \end{tabular}
  \end{center}

  \vspace{0.3em}
  Nenhum endereço começa com \texttt{/api}.\\
  Todos devolvem página, não JSON.
\end{frame}

\begin{frame}{As duas tabelas}
  \begin{columns}[T]
    \begin{column}{0.4\textwidth}
      \textbf{conta}
      \begin{itemize}
        \item id
        \item nome
        \item cpf
        \item senha
        \item saldo
      \end{itemize}
    \end{column}
    \begin{column}{0.55\textwidth}
      \textbf{movimentacao}
      \begin{itemize}
        \item id
        \item descricao --- ``Depósito'', ``PIX enviado para Bruno''
        \item valor --- negativo quando sai dinheiro
        \item data\_hora
        \item conta\_id --- de quem é
      \end{itemize}
    \end{column}
  \end{columns}

  \vspace{0.6em}
  Uma conta tem várias movimentações.\\
  É o 1{:}N da aula 05: o \texttt{conta\_id} faz o papel do \texttt{categoriaId}.
\end{frame}

\begin{frame}{Sessão: o crachá}
  Depois do login, como o Spring sabe quem está usando?

  \vspace{0.5em}
  Pensem num prédio com portaria:
  \begin{itemize}
    \item Na portaria, você mostra o documento (\textbf{login})
    \item Recebe um crachá (\textbf{sessão})
    \item Em cada andar, o segurança olha o crachá (\textbf{/conta})
    \item Na saída, devolve o crachá (\textbf{/sair})
  \end{itemize}

  \vspace{0.5em}
  No Java o crachá é o \texttt{HttpSession}.\\
  Nele guardamos o id da conta que entrou.
\end{frame}

\begin{frame}{As classes do projeto}
  \begin{center}
  \begin{tabular}{lll}
    \toprule
    Pasta & Arquivo & Para que serve \\
    \midrule
    \texttt{model} & \texttt{Conta} & a tabela conta \\
    \texttt{model} & \texttt{Movimentacao} & a tabela movimentacao \\
    \texttt{repository} & \texttt{ContaRepository} & fala com a tabela conta \\
    \texttt{repository} & \texttt{MovimentacaoRepository} & fala com a movimentacao \\
    \texttt{controller} & \texttt{AcessoController} & cadastro, login, sair \\
    \texttt{controller} & \texttt{ContaController} & saldo e operações \\
    \texttt{templates} & 3 arquivos \texttt{.html} & as três telas \\
    \bottomrule
  \end{tabular}
  \end{center}

  \vspace{0.5em}
  O CSS já está pronto. O resto a gente escreve junto.
\end{frame}

\begin{frame}{O que fica para cada dia}
  \begin{center}
  \begin{tabular}{ll}
    \toprule
    Quando & O que fica pronto \\
    \midrule
    Hoje, depois do intervalo & configuração, \texttt{Conta}, cadastro \\
    Amanhã, primeira parte & login, página da conta, depósito, saque \\
    Amanhã, segunda parte & PIX, \texttt{Movimentacao}, extrato \\
    \bottomrule
  \end{tabular}
  \end{center}

  \vspace{0.6em}
  Ao fim de hoje, cada cadastro já aparece no phpMyAdmin.
\end{frame}

% #############################################################################
\section{Configurando o projeto}

\begin{frame}{Passo 1 --- Abrir o projeto}
  \onde{Projeto \textbf{banco}}

  \begin{enumerate}
    \item Atualizem o repositório (\texttt{git pull})
    \item No VS Code: \textbf{File > Open Folder}
    \item Escolham a pasta \texttt{banco}
  \end{enumerate}

  \vspace{0.5em}
  O projeto foi gerado no \texttt{start.spring.io} com:
  \begin{itemize}
    \item Spring Web
    \item Thymeleaf
    \item Spring Data JPA
    \item MySQL Driver
  \end{itemize}
\end{frame}

\begin{frame}[fragile]{O que já veio pronto}
\begin{lstlisting}[basicstyle=\ttfamily\footnotesize]
banco/
  pom.xml
  src/main/java/com/entra21/banco/
    BancoApplication.java
  src/main/resources/
    application.properties
    static/css/estilo.css
\end{lstlisting}

  \vspace{0.4em}
  Só isso. Nenhum Controller, nenhum HTML.

  \vspace{0.4em}
  \begin{alertblock}{Não subam o Spring ainda}
    Com o JPA no projeto, o Spring quer um banco configurado.\\
    Se subir agora, dá erro. Primeiro o passo 2.
  \end{alertblock}
\end{frame}

\begin{frame}[fragile]{Onde cada arquivo vai ficar}
\begin{lstlisting}[basicstyle=\ttfamily\scriptsize]
banco/
  src/main/java/com/entra21/banco/
    BancoApplication.java
    model/Conta.java                          hoje
    model/Movimentacao.java                   amanha
    repository/ContaRepository.java           hoje
    repository/MovimentacaoRepository.java    amanha
    controller/AcessoController.java          hoje e amanha
    controller/ContaController.java           amanha
  src/main/resources/
    application.properties                    hoje
    static/css/estilo.css                     pronto
    templates/cadastro.html                   hoje
    templates/login.html                      amanha
    templates/conta.html                      amanha
\end{lstlisting}
\end{frame}

\begin{frame}[fragile]{Passo 2 --- Configurar o banco}
  \onde{\texttt{src/main/resources/application.properties}}

\begin{lstlisting}[language=props,basicstyle=\ttfamily\footnotesize]
spring.application.name=banco

spring.datasource.url=jdbc:mysql://localhost:3306/banco
spring.datasource.username=root
spring.datasource.password=

spring.jpa.hibernate.ddl-auto=update
spring.jpa.show-sql=true

server.servlet.session.tracking-modes=cookie
\end{lstlisting}

  A primeira linha já estava lá. Acrescentem o resto.
\end{frame}

\begin{frame}{O que cada linha faz}
  \begin{description}
    \item[datasource.url] Onde está o banco.\\
      \texttt{localhost:3306} é o MySQL do XAMPP.
      \texttt{/banco} é o nome que criamos.
    \item[username e password] \texttt{root}, sem senha.
    \item[ddl-auto=update] Criar e atualizar as tabelas sozinho.
    \item[show-sql=true] Mostrar no terminal o SQL que o Spring escreve.
    \item[tracking-modes] Guardar o crachá só no navegador,\\
      sem aparecer um código estranho no endereço.
  \end{description}
\end{frame}

\begin{frame}[fragile]{Passo 3 --- A classe Conta}
  \onde{Pasta nova \texttt{model}, arquivo \texttt{Conta.java}}

\begin{lstlisting}[language=Java,basicstyle=\ttfamily\scriptsize]
@Entity
public class Conta {

  @Id
  @GeneratedValue(strategy = GenerationType.IDENTITY)
  private Long id;

  private String nome;

  @Column(unique = true)
  private String cpf;

  private String senha;

  private double saldo;
\end{lstlisting}
\end{frame}

\begin{frame}{Passo 3 --- Completar a Conta}
  \onde{Ainda em \texttt{model/Conta.java}}

  \begin{itemize}
    \item Construtor vazio: \texttt{public Conta() \{ \}}
    \item Gets e sets dos cinco atributos
  \end{itemize}

  \vspace{0.5em}
  Os imports vêm de \texttt{jakarta.persistence}:\\
  \texttt{Entity}, \texttt{Id}, \texttt{GeneratedValue},
  \texttt{GenerationType}, \texttt{Column}.

  \vspace{0.5em}
  \begin{alertblock}{Cuidado no import}
    Se o VS Code oferecer \texttt{javax.persistence}, não aceitem.\\
    O certo é \texttt{jakarta.persistence}.
  \end{alertblock}
\end{frame}

\begin{frame}{Passo 4 --- Subir o Spring}
  Confiram: o MySQL está verde no XAMPP?

  \vspace{0.5em}
  No terminal, dentro da pasta \texttt{banco}:

  \vspace{0.3em}
  \texttt{./mvnw spring-boot:run}\\
  {\small No Windows: \texttt{.\textbackslash mvnw.cmd spring-boot:run}}

  \vspace{0.5em}
  Na primeira vez demora: o Maven baixa o JPA e o Thymeleaf.

  \vspace{0.5em}
  \begin{alertblock}{Um Spring por vez}
    Se outro projeto ainda estiver na porta 8080, parem o outro antes.
  \end{alertblock}
\end{frame}

\begin{frame}[fragile]{O que aparece no terminal}
  No meio das mensagens, procurem por esta:

\begin{lstlisting}[basicstyle=\ttfamily\scriptsize]
Hibernate: create table conta (id bigint not null auto_increment,
  cpf varchar(255), nome varchar(255), saldo float(53) not null,
  senha varchar(255), primary key (id))
\end{lstlisting}

  \vspace{0.4em}
  É o SQL que o Spring escreveu a partir da classe \texttt{Conta}.\\
  A gente não escreveu nenhuma linha dele.

  \vspace{0.4em}
  Aparece por causa do \texttt{show-sql=true}.
\end{frame}

\begin{frame}{Passo 5 --- Ver a tabela no phpMyAdmin}
  \begin{enumerate}
    \item Voltem para \texttt{http://localhost/phpmyadmin}
    \item Cliquem no banco \texttt{banco}, na esquerda
    \item Apareceu a tabela \texttt{conta}
    \item Cliquem nela e depois em \textbf{Estrutura}
  \end{enumerate}

  \vspace{0.5em}
  Estão lá as cinco colunas: id, cpf, nome, saldo, senha.

  \vspace{0.5em}
  A tabela ainda está vazia.\\
  Quem vai preencher é a tela de cadastro.
\end{frame}

\begin{frame}[fragile]{Passo 6 --- O ContaRepository}
  \onde{Pasta nova \texttt{repository}, arquivo \texttt{ContaRepository.java}}

\begin{lstlisting}[language=Java,basicstyle=\ttfamily\footnotesize]
public interface ContaRepository
    extends JpaRepository<Conta, Long> {

  Conta findByCpf(String cpf);
}
\end{lstlisting}

  \vspace{0.4em}
  Atenção: é \texttt{interface}, não \texttt{class}.

  \vspace{0.4em}
  Imports: \texttt{JpaRepository} e \texttt{Conta}.
\end{frame}

% #############################################################################
\section{Cadastro de contas}

\begin{frame}{Onde fica o HTML agora}
  Os HTML do Thymeleaf ficam numa pasta nova:

  \vspace{0.4em}
  \begin{center}
    \texttt{src/main/resources/templates/}
  \end{center}

  \vspace{0.4em}
  Criem essa pasta, ao lado da \texttt{static}.

  \vspace{0.5em}
  \begin{alertblock}{Não confundam}
    \texttt{static/} --- arquivos que vão do jeito que estão (CSS)\\
    \texttt{templates/} --- páginas que o Thymeleaf preenche antes
  \end{alertblock}
\end{frame}

\begin{frame}[fragile]{Passo 7 --- cadastro.html, o topo}
  \onde{Arquivo novo \texttt{templates/cadastro.html}}

\begin{lstlisting}[language=HTML,basicstyle=\ttfamily\scriptsize]
<!DOCTYPE html>
<html lang="pt-BR" xmlns:th="http://www.thymeleaf.org">
<head>
  <meta charset="UTF-8">
  <meta name="viewport" content="width=device-width, initial-scale=1.0">
  <title>Abrir conta</title>
  <link rel="stylesheet" href="/css/estilo.css">
</head>
<body>
  <div class="cartao">
    <p class="marca">Banco Entra21</p>
    <h1 class="titulo">Abrir conta</h1>

    <p class="alerta-sucesso" th:if="${sucesso}" th:text="${sucesso}"></p>
    <p class="alerta-erro" th:if="${erro}" th:text="${erro}"></p>
\end{lstlisting}

  O arquivo continua no próximo slide.
\end{frame}

\begin{frame}[fragile]{Passo 7 --- cadastro.html, o formulário}
  \onde{Continuação de \texttt{templates/cadastro.html}}

\begin{lstlisting}[language=HTML,basicstyle=\ttfamily\scriptsize]
    <form th:action="@{/cadastro}" method="post">
      <label>Nome completo</label>
      <input type="text" name="nome" required>

      <label>CPF (só números)</label>
      <input type="text" name="cpf" maxlength="11" required>

      <label>Senha</label>
      <input type="password" name="senha" required>

      <button class="botao" type="submit">Cadastrar</button>
    </form>

    <a class="link" th:href="@{/login}">Já tenho conta</a>
  </div>
</body>
</html>
\end{lstlisting}
\end{frame}

\begin{frame}{Lendo o cadastro.html}
  \begin{itemize}
    \item As classes (\texttt{cartao}, \texttt{botao}\ldots) já estão no CSS
    \item Os dois \texttt{th:if} mostram mensagem só quando existe
    \item O \texttt{th:action} manda o formulário para POST \texttt{/cadastro}
    \item Os \texttt{name} são \texttt{nome}, \texttt{cpf} e \texttt{senha}:\\
          iguais aos atributos da classe \texttt{Conta}
    \item \texttt{maxlength="11"}: o CPF só com números tem 11 dígitos
    \item \texttt{required}: o navegador não deixa enviar vazio
  \end{itemize}
\end{frame}

\begin{frame}[fragile]{Passo 8 --- AcessoController, mostrar a tela}
  \onde{Pasta nova \texttt{controller}, arquivo \texttt{AcessoController.java}}

\begin{lstlisting}[language=Java,basicstyle=\ttfamily\footnotesize]
@Controller
public class AcessoController {

  @Autowired
  private ContaRepository contaRepository;

  @GetMapping("/cadastro")
  public String telaCadastro() {
    return "cadastro";
  }
}
\end{lstlisting}

  \texttt{@Controller}, \textbf{não} \texttt{@RestController}.\\
  O \texttt{"cadastro"} é o nome do arquivo, sem o \texttt{.html}.
\end{frame}

\begin{frame}{Testando a tela}
  Parem o Spring (\texttt{Ctrl+C}) e subam de novo.

  \vspace{0.5em}
  Abram: \texttt{http://localhost:8080/cadastro}

  \vspace{0.5em}
  A tela de cadastro aparece, com o CSS.

  \vspace{0.5em}
  Se clicarem em Cadastrar agora, dá erro.\\
  Ainda não existe o POST. É o próximo passo.

  \vspace{0.5em}
  \begin{block}{Se aparecer só a palavra ``cadastro''}
    A classe está com \texttt{@RestController}. Troquem para \texttt{@Controller}.
  \end{block}
\end{frame}

\begin{frame}[fragile]{Passo 9 --- Receber o formulário}
  \onde{Ainda em \texttt{AcessoController.java}}

\begin{lstlisting}[language=Java,basicstyle=\ttfamily\scriptsize]
@PostMapping("/cadastro")
public String cadastrar(Conta conta, Model model,
                        RedirectAttributes redirect) {
  if (contaRepository.findByCpf(conta.getCpf()) != null) {
    model.addAttribute("erro", "Já existe uma conta com esse CPF");
    return "cadastro";
  }

  conta.setSaldo(0);
  contaRepository.save(conta);

  redirect.addFlashAttribute("sucesso", "Conta criada! Agora é só entrar.");
  return "redirect:/cadastro";
}
\end{lstlisting}

  Amanhã, quando o login existir, o último \texttt{return} muda.
\end{frame}

\begin{frame}{Lendo o POST do cadastro}
  \begin{description}
    \item[Conta conta] O Spring pega os \texttt{name} do formulário\\
      e preenche uma \texttt{Conta} nova: nome, cpf, senha.
    \item[findByCpf] Se já existe alguém com esse CPF, mostra o erro\\
      e devolve a mesma tela.
    \item[save] Grava no banco. Aqui acontece o \texttt{INSERT}.
    \item[addFlashAttribute] Um recado que sobrevive ao redirect.\\
      Aparece uma vez e some.
    \item[redirect] Depois do POST, sempre redirect.
  \end{description}
\end{frame}

\begin{frame}{Imports do AcessoController}
  \begin{itemize}
    \item \texttt{Conta} e \texttt{ContaRepository}
    \item \texttt{Controller} --- de \texttt{org.springframework.stereotype}
    \item \texttt{Autowired}, \texttt{GetMapping}, \texttt{PostMapping}
    \item \texttt{Model} --- de \texttt{org.springframework.ui}
    \item \texttt{RedirectAttributes}
  \end{itemize}

  \vspace{0.5em}
  Se o VS Code oferecer mais de um \texttt{Model},\\
  escolham o de \texttt{org.springframework.ui}.
\end{frame}

\begin{frame}{Passo 10 --- Cadastrar e conferir no banco}
  \begin{enumerate}
    \item Reiniciem o Spring
    \item Em \texttt{/cadastro}, criem uma conta com o nome de vocês
    \item A mensagem verde aparece
    \item No terminal, aparece um \texttt{insert into conta}
    \item No phpMyAdmin: tabela \texttt{conta}, aba \textbf{Visualizar}
  \end{enumerate}

  \vspace{0.5em}
  A linha está lá, com id 1 e saldo 0.

  \vspace{0.3em}
  Tentem cadastrar o mesmo CPF de novo: a mensagem vermelha aparece.
\end{frame}

\begin{frame}{A prova de que ficou gravado}
  Agora o teste que o \texttt{ArrayList} nunca passou:

  \vspace{0.5em}
  \begin{enumerate}
    \item Parem o Spring (\texttt{Ctrl+C})
    \item Subam de novo
    \item Olhem o phpMyAdmin
  \end{enumerate}

  \vspace{0.5em}
  A conta continua lá.\\
  Podem até desligar o computador. Amanhã ela ainda vai estar lá.
\end{frame}

\begin{frame}[fragile]{Olhando pela aba SQL}
  No phpMyAdmin, com o banco \texttt{banco} aberto, aba \textbf{SQL}:

\begin{lstlisting}[language=SQL,basicstyle=\ttfamily\footnotesize]
SELECT * FROM conta;

SELECT nome, cpf FROM conta WHERE nome LIKE 'A%';
\end{lstlisting}

  \vspace{0.4em}
  O primeiro mostra todas as contas.\\
  O segundo, só nome e CPF de quem começa com A.

  \vspace{0.4em}
  Não precisa decorar. É só para ver que o dado está no banco.
\end{frame}

\begin{frame}{Conferindo o Dia 1}
  \begin{itemize}
    \item [ ] XAMPP com MySQL verde
    \item [ ] O Spring sobe sem erro
    \item [ ] A tabela \texttt{conta} apareceu sozinha
    \item [ ] \texttt{/cadastro} abre com o CSS
    \item [ ] Cadastrar mostra a mensagem verde
    \item [ ] A conta aparece no phpMyAdmin
    \item [ ] CPF repetido mostra a mensagem vermelha
    \item [ ] Reiniciar o Spring não apaga a conta
  \end{itemize}
\end{frame}

\begin{frame}{Erros que mais aparecem}
  \begin{itemize}
    \item \texttt{Communications link failure} --- o MySQL do XAMPP
          está desligado
    \item \texttt{Unknown database 'banco'} --- faltou criar o banco
          no phpMyAdmin
    \item \texttt{Access denied for user 'root'} --- a senha no
          \texttt{application.properties} está errada
    \item Página de erro com \texttt{template might not exist} ---
          o arquivo não está em \texttt{templates/}, ou o nome no
          \texttt{return} está diferente
    \item Tabela não aparece --- faltou o \texttt{@Entity}
    \item Página sem CSS --- o \texttt{link} do CSS está diferente
          de \texttt{/css/estilo.css}
  \end{itemize}
\end{frame}

\begin{frame}{Fim do Dia 1}
  \begin{center}
    \Large Amanhã: login, saldo, depósito, saque,\\
    PIX e extrato.\\[0.8em]
    \normalsize
    Não apaguem o banco nem as contas.\\
    Amanhã a gente entra com elas.
  \end{center}
\end{frame}

% #############################################################################
\section{Dia 2 --- Login e sessão}

\begin{frame}{Onde paramos ontem}
  \begin{itemize}
    \item O XAMPP liga o MySQL
    \item O Spring cria a tabela \texttt{conta} sozinho
    \item \texttt{/cadastro} grava uma conta nova
    \item Os dados ficam no banco mesmo depois de reiniciar
  \end{itemize}

  \vspace{0.5em}
  \begin{block}{Antes de tudo, hoje}
    \begin{enumerate}
      \item XAMPP: Start no Apache e no MySQL
      \item Abram o projeto \texttt{banco} e subam o Spring
      \item Confiram no phpMyAdmin: as contas de ontem estão lá?
    \end{enumerate}
  \end{block}
\end{frame}

\begin{frame}{O plano de hoje}
  \begin{block}{Parte 3 --- 18h30 às 20h}
    Tela de login, sessão, página da conta, depósito e saque.
  \end{block}

  \vspace{0.4em}
  \begin{block}{Parte 4 --- 20h20 às 22h}
    PIX, extrato e o teste final com pessoas do gerador.
  \end{block}

  \vspace{0.6em}
  Ao final, cada um terá um banco funcionando,\\
  com o dinheiro indo de uma conta para outra.
\end{frame}

\begin{frame}[fragile]{Passo 11 --- login.html, o topo}
  \onde{Arquivo novo \texttt{templates/login.html}.\\
        Copiem o \texttt{<head>} do cadastro e troquem o título para Entrar.}

\begin{lstlisting}[language=HTML,basicstyle=\ttfamily\scriptsize]
<body>
  <div class="cartao">
    <p class="marca">Banco Entra21</p>
    <h1 class="titulo">Entrar</h1>

    <p class="alerta-sucesso" th:if="${sucesso}" th:text="${sucesso}"></p>
    <p class="alerta-erro" th:if="${erro}" th:text="${erro}"></p>
\end{lstlisting}

  É o mesmo começo do cadastro, só muda o título.
\end{frame}

\begin{frame}[fragile]{Passo 11 --- login.html, o formulário}
  \onde{Continuação de \texttt{templates/login.html}}

\begin{lstlisting}[language=HTML,basicstyle=\ttfamily\scriptsize]
    <form th:action="@{/login}" method="post">
      <label>CPF</label>
      <input type="text" name="cpf" maxlength="11" required>

      <label>Senha</label>
      <input type="password" name="senha" required>

      <button class="botao" type="submit">Entrar</button>
    </form>

    <a class="link" th:href="@{/cadastro}">Abrir uma conta</a>
  </div>
</body>
</html>
\end{lstlisting}
\end{frame}

\begin{frame}[fragile]{Passo 12 --- Mostrar o login}
  \onde{\texttt{AcessoController.java}, acrescentar}

\begin{lstlisting}[language=Java,basicstyle=\ttfamily\footnotesize]
@GetMapping("/")
public String inicio() {
  return "redirect:/login";
}

@GetMapping("/login")
public String telaLogin() {
  return "login";
}
\end{lstlisting}

  \vspace{0.3em}
  Quem abrir \texttt{http://localhost:8080} cai direto no login.
\end{frame}

\begin{frame}[fragile]{Passo 13 --- Cadastro agora leva ao login}
  \onde{\texttt{AcessoController.java}, no fim do POST \texttt{/cadastro}}

  Ontem era:

\begin{lstlisting}[language=Java,basicstyle=\ttfamily\footnotesize]
  return "redirect:/cadastro";
\end{lstlisting}

  Troquem por:

\begin{lstlisting}[language=Java,basicstyle=\ttfamily\footnotesize]
  return "redirect:/login";
\end{lstlisting}

  \vspace{0.3em}
  A mensagem ``Conta criada!'' agora aparece na tela de login.\\
  Por isso o \texttt{login.html} também tem o \texttt{th:if} do sucesso.
\end{frame}

\begin{frame}{Como funciona o login}
  A pessoa manda CPF e senha. O servidor:

  \vspace{0.5em}
  \begin{enumerate}
    \item Procura a conta pelo CPF (\texttt{findByCpf})
    \item Se não achou, erro
    \item Se achou, mas a senha é diferente, erro
    \item Se deu tudo certo, guarda o id da conta na sessão
    \item Manda para \texttt{/conta}
  \end{enumerate}

  \vspace{0.5em}
  O erro é o mesmo nos dois casos: ``CPF ou senha incorretos''.\\
  Assim ninguém descobre quais CPFs têm conta.
\end{frame}

\begin{frame}[fragile]{Passo 14 --- POST do login}
  \onde{\texttt{AcessoController.java}, acrescentar}

\begin{lstlisting}[language=Java,basicstyle=\ttfamily\scriptsize]
@PostMapping("/login")
public String entrar(@RequestParam String cpf, @RequestParam String senha,
                     Model model, HttpSession session) {
  Conta conta = contaRepository.findByCpf(cpf);

  if (conta == null || !conta.getSenha().equals(senha)) {
    model.addAttribute("erro", "CPF ou senha incorretos");
    return "login";
  }

  session.setAttribute("contaId", conta.getId());
  return "redirect:/conta";
}
\end{lstlisting}

\end{frame}

\begin{frame}{Lendo o POST do login}
  \begin{description}
    \item[findByCpf] Procura a conta no banco. Se não achar, vem \texttt{null}.
    \item[||] Quer dizer ``ou''. Não achou a conta \textbf{ou} a senha é outra.
    \item[equals] Compara texto. Com \texttt{String}, nunca usem \texttt{==}.
    \item[!] Quer dizer ``não''. \texttt{!equals} é ``não é igual''.
    \item[return] Devolve a mesma tela de login, agora com o erro.
  \end{description}
\end{frame}

\begin{frame}[fragile]{A sessão por dentro}
\begin{lstlisting}[language=Java,basicstyle=\ttfamily\footnotesize]
session.setAttribute("contaId", conta.getId());
\end{lstlisting}

  \vspace{0.3em}
  É o crachá sendo entregue.\\
  Guardamos só o \textbf{id}, com a etiqueta \texttt{"contaId"}.

  \vspace{0.5em}
  Cada navegador tem a sua sessão.\\
  Se abrirem uma janela anônima, podem entrar com outra conta\\
  ao mesmo tempo.

  \vspace{0.5em}
  Import: \texttt{jakarta.servlet.http.HttpSession}.
\end{frame}

\begin{frame}[fragile]{Passo 15 --- Sair}
  \onde{\texttt{AcessoController.java}, acrescentar}

\begin{lstlisting}[language=Java,basicstyle=\ttfamily\footnotesize]
@GetMapping("/sair")
public String sair(HttpSession session) {
  session.invalidate();
  return "redirect:/login";
}
\end{lstlisting}

  \vspace{0.4em}
  \texttt{invalidate} rasga o crachá.\\
  Tudo que estava na sessão é esquecido.
\end{frame}

\begin{frame}{Uma palavra sobre a senha}
  Estamos guardando a senha do jeito que a pessoa digitou.\\
  Quem abrir o phpMyAdmin consegue ler.

  \vspace{0.6em}
  Isso serve só para aprender.

  \vspace{0.6em}
  Em sistema de verdade:
  \begin{itemize}
    \item a senha é guardada embaralhada, sem volta (\textit{hash})
    \item o login é feito pelo \textbf{Spring Security}
  \end{itemize}

  \vspace{0.6em}
  Nunca usem uma senha de verdade de vocês nesse projeto.
\end{frame}

% #############################################################################
\section{A página da conta}

\begin{frame}[fragile]{Passo 16 --- ContaController}
  \onde{Arquivo novo \texttt{controller/ContaController.java}}

\begin{lstlisting}[language=Java,basicstyle=\ttfamily\scriptsize]
@Controller
public class ContaController {

  @Autowired
  private ContaRepository contaRepository;

  public Conta contaLogada(HttpSession session) {
    Long id = (Long) session.getAttribute("contaId");
    if (id == null) {
      return null;
    }
    return contaRepository.findById(id).orElse(null);
  }
}
\end{lstlisting}

\end{frame}

\begin{frame}{Lendo o contaLogada}
  Uma função que olha o crachá e devolve a conta,
  ou \texttt{null} se ninguém entrou.

  \vspace{0.3em}
  \begin{description}
    \item[getAttribute] Pega o que está na sessão com a etiqueta \texttt{"contaId"}.
    \item[(Long)] A sessão guarda tudo como \texttt{Object}.\\
      O \texttt{(Long)} avisa ao Java: ``isso aqui é um Long''.
    \item[id == null] Ninguém fez login nesse navegador.
    \item[findById] Busca a conta no banco pelo id.
    \item[orElse(null)] Se não achar, devolve \texttt{null}.
  \end{description}

  \vspace{0.4em}
  Por que buscar no banco de novo?\\
  Porque o saldo pode ter mudado. Por exemplo, alguém mandou um PIX.
\end{frame}

\begin{frame}[fragile]{Passo 17 --- GET /conta}
  \onde{\texttt{ContaController.java}, acrescentar}

\begin{lstlisting}[language=Java,basicstyle=\ttfamily\footnotesize]
@GetMapping("/conta")
public String telaConta(HttpSession session, Model model) {
  Conta conta = contaLogada(session);
  if (conta == null) {
    return "redirect:/login";
  }

  model.addAttribute("conta", conta);
  return "conta";
}
\end{lstlisting}

  Sem crachá, volta para o login.\\
  Com crachá, coloca a conta na bandeja e mostra a página.
\end{frame}

\begin{frame}[fragile]{Passo 18 --- conta.html, o topo}
  \onde{Arquivo novo \texttt{templates/conta.html}.\\
        Copiem o \texttt{<head>} do cadastro e troquem o título para Minha conta.}

\begin{lstlisting}[language=HTML,basicstyle=\ttfamily\scriptsize]
<body>
  <div class="cartao cartao-larga">
    <div class="topo">
      <p class="marca">Banco Entra21</p>
      <a class="botao-secundario" th:href="@{/sair}">Sair</a>
    </div>

    <p class="ola">Olá, <strong th:text="${conta.nome}">Nome</strong></p>
\end{lstlisting}

  \texttt{\$\{conta.nome\}}: a conta que o Controller pôs na bandeja,\\
  e dela o \texttt{getNome()}.
\end{frame}

\begin{frame}[fragile]{Passo 18 --- conta.html, o saldo}
  \onde{Continuação de \texttt{templates/conta.html}}

\begin{lstlisting}[language=HTML,basicstyle=\ttfamily\scriptsize]
    <div class="saldo">
      <p class="saldo-rotulo">Saldo disponível</p>
      <p class="saldo-valor" th:text="'R$ ' +
        ${#numbers.formatDecimal(conta.saldo, 1, 'POINT', 2, 'COMMA')}">
        R$ 0,00</p>
    </div>

    <p class="alerta-sucesso" th:if="${sucesso}" th:text="${sucesso}"></p>
    <p class="alerta-erro" th:if="${erro}" th:text="${erro}"></p>
\end{lstlisting}

  A linha comprida do saldo é explicada no próximo slide.
\end{frame}

\begin{frame}[fragile]{Formatando o dinheiro}
\begin{lstlisting}[basicstyle=\ttfamily\footnotesize]
${#numbers.formatDecimal(conta.saldo, 1, 'POINT', 2, 'COMMA')}
\end{lstlisting}

  \vspace{0.3em}
  Sem isso, 1234.5 aparece como \texttt{1234.5}.\\
  Com isso, aparece \texttt{1.234,50}.

  \vspace{0.4em}
  \begin{center}
  \begin{tabular}{ll}
    \toprule
    Pedaço & Significa \\
    \midrule
    \texttt{conta.saldo} & o número \\
    \texttt{1} & pelo menos 1 dígito antes da vírgula \\
    \texttt{'POINT'} & separa os milhares com ponto \\
    \texttt{2} & duas casas depois da vírgula \\
    \texttt{'COMMA'} & os centavos vêm depois de uma vírgula \\
    \bottomrule
  \end{tabular}
  \end{center}
\end{frame}

\begin{frame}[fragile]{Passo 19 --- conta.html, o depósito}
  \onde{Continuação de \texttt{templates/conta.html}}

\begin{lstlisting}[language=HTML,basicstyle=\ttfamily\scriptsize]
    <div class="operacoes">
      <form class="operacao" th:action="@{/deposito}" method="post">
        <h3>Depósito</h3>
        <label>Valor</label>
        <input type="number" name="valor" step="0.01" min="0.01" required>
        <button class="botao" type="submit">Depositar</button>
      </form>
\end{lstlisting}

  \texttt{step="0.01"} deixa digitar centavos.\\
  \texttt{min="0.01"} não deixa digitar zero nem negativo.
\end{frame}

\begin{frame}[fragile]{Passo 19 --- conta.html, o saque}
  \onde{Continuação de \texttt{templates/conta.html}}

\begin{lstlisting}[language=HTML,basicstyle=\ttfamily\scriptsize]
      <form class="operacao" th:action="@{/saque}" method="post">
        <h3>Saque</h3>
        <label>Valor</label>
        <input type="number" name="valor" step="0.01" min="0.01" required>
        <button class="botao" type="submit">Sacar</button>
      </form>
    </div>
  </div>
</body>
</html>
\end{lstlisting}

  O primeiro \texttt{</div>} fecha as operações.\\
  O segundo fecha o cartão.
\end{frame}

\begin{frame}{Testando a página da conta}
  \begin{enumerate}
    \item Reiniciem o Spring
    \item Abram \texttt{http://localhost:8080}: cai no login
    \item Entrem com a conta de ontem
    \item Aparece o nome e \texttt{R\$ 0,00}
    \item Cliquem em Sair e tentem abrir \texttt{/conta} direto
    \item Volta para o login: sem crachá, não entra
  \end{enumerate}

  \vspace{0.4em}
  Os botões de depósito e saque ainda dão erro.\\
  É o próximo passo.
\end{frame}

% #############################################################################
\section{Depósito e saque}

\begin{frame}[fragile]{Passo 20 --- Depósito, as conferências}
  \onde{\texttt{ContaController.java}, acrescentar}

\begin{lstlisting}[language=Java,basicstyle=\ttfamily\scriptsize]
@PostMapping("/deposito")
public String depositar(@RequestParam double valor, HttpSession session,
                        RedirectAttributes redirect) {
  Conta conta = contaLogada(session);
  if (conta == null) {
    return "redirect:/login";
  }

  if (valor <= 0) {
    redirect.addFlashAttribute("erro", "Digite um valor maior que zero");
    return "redirect:/conta";
  }
\end{lstlisting}

  O método continua no próximo slide.
\end{frame}

\begin{frame}[fragile]{Passo 20 --- Depósito, somando}
  \onde{Ainda no método \texttt{depositar}}

\begin{lstlisting}[language=Java,basicstyle=\ttfamily\footnotesize]
  conta.setSaldo(conta.getSaldo() + valor);
  contaRepository.save(conta);

  redirect.addFlashAttribute("sucesso", "Depósito feito");
  return "redirect:/conta";
}
\end{lstlisting}
\end{frame}

\begin{frame}{Lendo o depósito}
  \begin{enumerate}
    \item \texttt{@RequestParam double valor} --- vem do \texttt{name="valor"}
    \item Confere o crachá
    \item Recusa valor zero ou negativo
    \item Soma no saldo do objeto
    \item \texttt{save} --- como a conta já tem id, vira \texttt{UPDATE}
    \item Deixa o recado verde e volta para \texttt{/conta}
  \end{enumerate}

  \vspace{0.5em}
  O \texttt{/conta} busca a conta de novo no banco.\\
  Por isso o saldo aparece atualizado.
\end{frame}

\begin{frame}[fragile]{Passo 21 --- Saque, as conferências}
  \onde{\texttt{ContaController.java}, acrescentar}

\begin{lstlisting}[language=Java,basicstyle=\ttfamily\scriptsize]
@PostMapping("/saque")
public String sacar(@RequestParam double valor, HttpSession session,
                    RedirectAttributes redirect) {
  Conta conta = contaLogada(session);
  if (conta == null) {
    return "redirect:/login";
  }

  if (valor <= 0) {
    redirect.addFlashAttribute("erro", "Digite um valor maior que zero");
    return "redirect:/conta";
  }
\end{lstlisting}

  Até aqui é igual ao depósito.
\end{frame}

\begin{frame}[fragile]{Passo 21 --- Saque, tirando do saldo}
  \onde{Ainda no método \texttt{sacar}}

\begin{lstlisting}[language=Java,basicstyle=\ttfamily\footnotesize]
  if (conta.getSaldo() < valor) {
    redirect.addFlashAttribute("erro", "Saldo insuficiente");
    return "redirect:/conta";
  }

  conta.setSaldo(conta.getSaldo() - valor);
  contaRepository.save(conta);

  redirect.addFlashAttribute("sucesso", "Saque feito");
  return "redirect:/conta";
}
\end{lstlisting}
\end{frame}

\begin{frame}{O saque é o depósito com uma conferência a mais}
  \begin{center}
  \begin{tabular}{ll}
    \toprule
    Depósito & Saque \\
    \midrule
    confere o crachá & confere o crachá \\
    valor maior que zero & valor maior que zero \\
    --- & \textbf{saldo suficiente?} \\
    saldo \textbf{+} valor & saldo \textbf{-} valor \\
    \texttt{save} & \texttt{save} \\
    \bottomrule
  \end{tabular}
  \end{center}

  \vspace{0.6em}
  Igual à venda da aula 06:\\
  primeiro confere o estoque, depois tira.
\end{frame}

\begin{frame}{Imports do ContaController}
  \begin{itemize}
    \item \texttt{Conta} e \texttt{ContaRepository}
    \item \texttt{Controller}, \texttt{Autowired}
    \item \texttt{GetMapping}, \texttt{PostMapping}, \texttt{RequestParam}
    \item \texttt{Model} --- de \texttt{org.springframework.ui}
    \item \texttt{RedirectAttributes}
    \item \texttt{HttpSession} --- de \texttt{jakarta.servlet.http}
  \end{itemize}

  \vspace{0.5em}
  Se faltar um import, o VS Code sublinha em vermelho.\\
  Aceitem o import que ele oferece, olhando o nome do pacote.
\end{frame}

\begin{frame}{Conferindo a Parte 3}
  \begin{itemize}
    \item [ ] \texttt{/} leva para o login
    \item [ ] Senha errada mostra ``CPF ou senha incorretos''
    \item [ ] Login certo mostra o nome e o saldo
    \item [ ] Depositar 500 mostra \texttt{R\$ 500,00}
    \item [ ] Sacar mais do que tem mostra ``Saldo insuficiente''
    \item [ ] No phpMyAdmin o saldo mudou
    \item [ ] Sair e abrir \texttt{/conta} volta para o login
    \item [ ] F5 depois do depósito \textbf{não} deposita de novo
  \end{itemize}
\end{frame}

\begin{frame}{Erros que mais aparecem}
  \begin{itemize}
    \item \texttt{Required parameter 'valor' is not present} ---
          o \texttt{name} do input não é \texttt{valor}
    \item Saldo não muda --- faltou o \texttt{contaRepository.save(conta)}
    \item \texttt{Exception evaluating SpringEL} no \texttt{conta.html} ---
          o Controller não fez \texttt{model.addAttribute("conta", ...)}
    \item Sempre volta para o login --- no \texttt{setAttribute} e no
          \texttt{getAttribute} a etiqueta está diferente
    \item Mensagem não aparece --- usaram \texttt{model} em vez de
          \texttt{redirect.addFlashAttribute} antes de um redirect
  \end{itemize}
\end{frame}

\begin{frame}{Intervalo}
  \begin{center}
    \Large Retomamos às \textbf{20h20}\\[0.8em]
    \normalsize
    Na volta: PIX entre contas e o extrato.\\
    Deixem pelo menos duas contas cadastradas.
  \end{center}
\end{frame}

% #############################################################################
\section{PIX}

\begin{frame}{O que é um PIX no nosso banco}
  Dinheiro sai de uma conta e entra em outra.\\
  Quem recebe é encontrado pelo CPF.

  \vspace{0.5em}
  O servidor faz nesta ordem:
  \begin{enumerate}
    \item Confere o crachá: quem está mandando
    \item Valor maior que zero
    \item Acha a conta de destino pelo CPF
    \item Não deixa mandar para si mesmo
    \item Confere se quem manda tem saldo
    \item Tira de um, põe no outro, grava os dois
  \end{enumerate}

  \vspace{0.3em}
  Primeiro confere \textbf{tudo}. Só depois mexe no dinheiro.
\end{frame}

\begin{frame}[fragile]{Passo 22 --- O formulário do PIX}
  \onde{\texttt{templates/conta.html}, dentro da \texttt{div operacoes},\\
        depois do formulário de saque}

\begin{lstlisting}[language=HTML,basicstyle=\ttfamily\scriptsize]
<form class="operacao" th:action="@{/pix}" method="post">
  <h3>PIX</h3>
  <label>CPF de quem recebe</label>
  <input type="text" name="cpfDestino" maxlength="11" required>
  <label>Valor</label>
  <input type="number" name="valor" step="0.01" min="0.01" required>
  <button class="botao" type="submit">Enviar PIX</button>
</form>
\end{lstlisting}

  Dois campos: \texttt{cpfDestino} e \texttt{valor}.\\
  O CSS já coloca os três formulários lado a lado.
\end{frame}

\begin{frame}[fragile]{Passo 23 --- PIX, o começo}
  \onde{\texttt{ContaController.java}, acrescentar}

\begin{lstlisting}[language=Java,basicstyle=\ttfamily\scriptsize]
@Transactional
@PostMapping("/pix")
public String pix(@RequestParam String cpfDestino, @RequestParam double valor,
                  HttpSession session, RedirectAttributes redirect) {
  Conta origem = contaLogada(session);
  if (origem == null) {
    return "redirect:/login";
  }

  if (valor <= 0) {
    redirect.addFlashAttribute("erro", "Digite um valor maior que zero");
    return "redirect:/conta";
  }
\end{lstlisting}

  \texttt{origem} é quem manda. Quem recebe vem no próximo slide.
\end{frame}

\begin{frame}[fragile]{Passo 23 --- PIX, achando quem recebe}
  \onde{Ainda no método \texttt{pix}}

\begin{lstlisting}[language=Java,basicstyle=\ttfamily\scriptsize]
  Conta destino = contaRepository.findByCpf(cpfDestino);
  if (destino == null) {
    redirect.addFlashAttribute("erro", "Nenhuma conta com esse CPF");
    return "redirect:/conta";
  }

  if (destino.getId().equals(origem.getId())) {
    redirect.addFlashAttribute("erro", "Não dá para mandar PIX para você mesmo");
    return "redirect:/conta";
  }
\end{lstlisting}

  Por que \texttt{equals} para comparar os ids?\\
  \texttt{Long} com L maiúsculo é objeto. Objeto se compara com \texttt{equals}.
\end{frame}

\begin{frame}[fragile]{Passo 23 --- PIX, mexendo no dinheiro}
  \onde{Ainda no método \texttt{pix}}

\begin{lstlisting}[language=Java,basicstyle=\ttfamily\scriptsize]
  if (origem.getSaldo() < valor) {
    redirect.addFlashAttribute("erro", "Saldo insuficiente");
    return "redirect:/conta";
  }

  origem.setSaldo(origem.getSaldo() - valor);
  destino.setSaldo(destino.getSaldo() + valor);
  contaRepository.save(origem);
  contaRepository.save(destino);

  redirect.addFlashAttribute("sucesso",
      "PIX enviado para " + destino.getNome());
  return "redirect:/conta";
}
\end{lstlisting}

  Tira de um, põe no outro, grava os dois.
\end{frame}

\begin{frame}{@Transactional: tudo ou nada}
  Imaginem que a luz acaba bem no meio do PIX:

  \vspace{0.4em}
  \begin{itemize}
    \item O \texttt{save(origem)} já rodou: a Ana perdeu 100
    \item O \texttt{save(destino)} não rodou: o Bruno não recebeu
    \item Os 100 reais sumiram
  \end{itemize}

  \vspace{0.5em}
  O \texttt{@Transactional} em cima do método diz ao banco:\\
  ``ou grava tudo que tem aqui dentro, ou não grava nada''.

  \vspace{0.5em}
  \begin{alertblock}{Cuidado no import}
    O certo é \texttt{org.springframework.transaction.annotation.Transactional}.
  \end{alertblock}
\end{frame}

\begin{frame}{Conferindo o PIX}
  Usem duas contas: uma no navegador normal,\\
  outra numa janela anônima.

  \vspace{0.5em}
  \begin{itemize}
    \item [ ] PIX para um CPF que não existe: mensagem vermelha
    \item [ ] PIX para o próprio CPF: mensagem vermelha
    \item [ ] PIX maior que o saldo: ``Saldo insuficiente''
    \item [ ] PIX certo: o saldo cai de um lado
    \item [ ] F5 na janela anônima: o saldo sobe do outro lado
  \end{itemize}
\end{frame}

% #############################################################################
\section{Extrato}

\begin{frame}{Por que outra tabela}
  O saldo diz \textbf{quanto} tem.\\
  Não diz \textbf{o que aconteceu}.

  \vspace{0.5em}
  Para o extrato, cada operação vira uma linha nova:

  \vspace{0.4em}
  \begin{center}
  \small
  \begin{tabular}{lllll}
    \toprule
    id & conta\_id & descricao & valor & data\_hora \\
    \midrule
    1 & 1 & Depósito & 500.50 & 23/09 19:02 \\
    2 & 1 & Saque & -100.00 & 23/09 19:05 \\
    3 & 1 & PIX enviado para Bruno Lima & -150.25 & 23/09 20:31 \\
    4 & 2 & PIX recebido de Ana Souza & 150.25 & 23/09 20:31 \\
    \bottomrule
  \end{tabular}
  \end{center}

  \vspace{0.4em}
  O \texttt{conta\_id} diz de quem é a linha.\\
  Um PIX gera \textbf{duas} linhas: uma em cada conta.
\end{frame}

\begin{frame}{Saiu ou entrou?}
  Na coluna \texttt{valor}:

  \vspace{0.4em}
  \begin{itemize}
    \item Entrou dinheiro: número \textbf{positivo}\\
          depósito, PIX recebido
    \item Saiu dinheiro: número \textbf{negativo}\\
          saque, PIX enviado
  \end{itemize}

  \vspace{0.5em}
  Assim, no extrato, é fácil pintar:\\
  verde para positivo, vermelho para negativo.
\end{frame}

\begin{frame}[fragile]{Passo 24 --- A classe Movimentacao}
  \onde{Arquivo novo \texttt{model/Movimentacao.java}}

\begin{lstlisting}[language=Java,basicstyle=\ttfamily\scriptsize]
@Entity
public class Movimentacao {

  @Id
  @GeneratedValue(strategy = GenerationType.IDENTITY)
  private Long id;

  private String descricao;

  private double valor;

  private LocalDateTime dataHora;

  @ManyToOne
  private Conta conta;
\end{lstlisting}

  Construtor vazio e gets e sets dos cinco atributos.
\end{frame}

\begin{frame}{@ManyToOne --- muitas para uma}
  \begin{center}
  \begin{tikzpicture}[
      box/.style={draw, rounded corners, align=center,
                  minimum width=2.4cm, minimum height=0.8cm,
                  font=\small},
      meio/.style={draw, rounded corners, align=center,
                   minimum width=2.4cm, minimum height=0.55cm,
                   font=\scriptsize},
      seta/.style={-{Latex}, thick}
    ]
    \node[box, fill=blue!15] (c) {Conta da Ana};
    \node[meio, fill=yellow!25, right=2.2cm of c, yshift=0.8cm] (m1) {Depósito};
    \node[meio, fill=yellow!25, right=2.2cm of c] (m2) {Saque};
    \node[meio, fill=yellow!25, right=2.2cm of c, yshift=-0.8cm] (m3) {PIX enviado};

    \draw[seta] (m1.west) -- (c);
    \draw[seta] (m2.west) -- (c);
    \draw[seta] (m3.west) -- (c);
  \end{tikzpicture}
  \end{center}

  \vspace{0.4em}
  \textbf{Muitas} movimentações apontam para \textbf{uma} conta.\\
  No banco, isso vira a coluna \texttt{conta\_id}.

  \vspace{0.3em}
  Na aula 05 a gente escrevia \texttt{int categoriaId} na mão.\\
  Agora a gente põe a \texttt{Conta} inteira e o JPA guarda só o id.
\end{frame}

\begin{frame}{LocalDateTime}
  \texttt{LocalDateTime} guarda dia e hora juntos.

  \vspace{0.5em}
  \texttt{LocalDateTime.now()} pega o momento de agora.

  \vspace{0.5em}
  No banco vira a coluna \texttt{data\_hora}.\\
  O JPA troca o \texttt{dataHora} do Java pelo \texttt{data\_hora} do banco.

  \vspace{0.5em}
  Imports:
  \begin{itemize}
    \item \texttt{java.time.LocalDateTime}
    \item \texttt{ManyToOne} e as outras anotações de \texttt{jakarta.persistence}
  \end{itemize}
\end{frame}

\begin{frame}[fragile]{Passo 25 --- MovimentacaoRepository}
  \onde{Arquivo novo \texttt{repository/MovimentacaoRepository.java}}

\begin{lstlisting}[language=Java,basicstyle=\ttfamily\footnotesize]
public interface MovimentacaoRepository
    extends JpaRepository<Movimentacao, Long> {

  List<Movimentacao> findByContaOrderByDataHoraDesc(Conta conta);
}
\end{lstlisting}

  \vspace{0.2em}
  \begin{center}
  \small
  \begin{tabular}{ll}
    \toprule
    Pedaço & Significa \\
    \midrule
    \texttt{findBy} & buscar onde \\
    \texttt{Conta} & a conta é esta \\
    \texttt{OrderBy} & em ordem de \\
    \texttt{DataHora} & data e hora \\
    \texttt{Desc} & da mais nova para a mais velha \\
    \bottomrule
  \end{tabular}
  \end{center}
\end{frame}

\begin{frame}[fragile]{Passo 26 --- Uma função para registrar}
  \onde{\texttt{ContaController.java}, acrescentar}

\begin{lstlisting}[language=Java,basicstyle=\ttfamily\footnotesize]
@Autowired
private MovimentacaoRepository movimentacaoRepository;

public void registrar(Conta conta, String descricao, double valor) {
  Movimentacao mov = new Movimentacao();
  mov.setConta(conta);
  mov.setDescricao(descricao);
  mov.setValor(valor);
  mov.setDataHora(LocalDateTime.now());
  movimentacaoRepository.save(mov);
}
\end{lstlisting}

  Em vez de repetir essas linhas em cada operação,\\
  chamamos \texttt{registrar} com os três dados.
\end{frame}

\begin{frame}[fragile]{Passo 27 --- Registrar cada operação}
  \onde{\texttt{ContaController.java}, logo depois de cada \texttt{save}}

  No depósito:
\begin{lstlisting}[language=Java,basicstyle=\ttfamily\footnotesize]
registrar(conta, "Depósito", valor);
\end{lstlisting}

  No saque, com o menos na frente:
\begin{lstlisting}[language=Java,basicstyle=\ttfamily\footnotesize]
registrar(conta, "Saque", -valor);
\end{lstlisting}

  No PIX, uma linha para cada conta:
\begin{lstlisting}[language=Java,basicstyle=\ttfamily\footnotesize]
registrar(origem, "PIX enviado para " + destino.getNome(), -valor);
registrar(destino, "PIX recebido de " + origem.getNome(), valor);
\end{lstlisting}
\end{frame}

\begin{frame}[fragile]{Passo 28 --- Mandar o extrato para a página}
  \onde{\texttt{ContaController.java}, no GET \texttt{/conta}}

\begin{lstlisting}[language=Java,basicstyle=\ttfamily\footnotesize]
model.addAttribute("conta", conta);
model.addAttribute("movimentacoes",
    movimentacaoRepository.findByContaOrderByDataHoraDesc(conta));
return "conta";
\end{lstlisting}

  \vspace{0.4em}
  A linha do meio é nova.\\
  Coloca na bandeja a lista de movimentações dessa conta.

  \vspace{0.4em}
  Imports novos no Controller:\\
  \texttt{Movimentacao}, \texttt{MovimentacaoRepository},
  \texttt{LocalDateTime}, \texttt{Transactional}.
\end{frame}

\begin{frame}[fragile]{Passo 29 --- O extrato, o cabeçalho}
  \onde{\texttt{templates/conta.html}, depois da \texttt{div operacoes},\\
        antes de fechar a \texttt{div cartao}}

\begin{lstlisting}[language=HTML,basicstyle=\ttfamily\footnotesize]
<h2 class="subtitulo">Extrato</h2>

<table class="extrato">
  <tr>
    <th>Data</th>
    <th>Descrição</th>
    <th>Valor</th>
  </tr>
\end{lstlisting}

  Até aqui é HTML comum: o título e a primeira linha da tabela.
\end{frame}

\begin{frame}[fragile]{Passo 29 --- O extrato, uma linha por movimentação}
  \onde{Continuação, logo abaixo}

\begin{lstlisting}[language=HTML,basicstyle=\ttfamily\scriptsize]
  <tr th:each="mov : ${movimentacoes}">
    <td th:text="${#temporals.format(mov.dataHora, 'dd/MM/yyyy HH:mm')}"></td>
    <td th:text="${mov.descricao}"></td>
    <td th:text="${#numbers.formatDecimal(mov.valor, 1, 'POINT', 2, 'COMMA')}"
        th:classappend="${mov.valor < 0} ? 'saida' : 'entrada'"></td>
  </tr>
</table>

<p class="vazio" th:if="${movimentacoes.isEmpty()}">Nenhuma movimentação ainda.</p>
\end{lstlisting}
\end{frame}

\begin{frame}{Lendo o extrato}
  \begin{description}
    \item[th:each] Repete o \texttt{<tr>} para cada \texttt{mov} da lista.
    \item[\#temporals.format] Escreve a data do jeito brasileiro:\\
      \texttt{23/09/2026 20:31}.
    \item[formatDecimal] O mesmo que usamos no saldo.
    \item[th:classappend] Acrescenta uma classe do CSS:\\
      \texttt{saida} (vermelho) se o valor for negativo,\\
      \texttt{entrada} (verde) se não for.
    \item[isEmpty()] Se a lista estiver vazia, mostra o aviso.
  \end{description}

  \vspace{0.3em}
  O \texttt{? :} é um \texttt{if} curto: ``é negativo? então saida, senão entrada''.
\end{frame}

\begin{frame}{Testando o extrato}
  \begin{enumerate}
    \item Reiniciem o Spring
    \item No terminal: \texttt{create table movimentacao}
    \item No phpMyAdmin: a tabela nova, com a coluna \texttt{conta\_id}
    \item Entrem, depositem, saquem e mandem um PIX
    \item Cada operação aparece no extrato, a mais nova em cima
    \item Entrem na outra conta: o PIX recebido aparece em verde
  \end{enumerate}

  \vspace{0.4em}
  O que foi feito antes do extrato não aparece.\\
  Só a partir de agora cada operação é registrada.
\end{frame}

% #############################################################################
\section{Teste final}

\begin{frame}{Pessoas de mentira, CPF válido}
  Para testar, vamos criar pessoas que não existem.

  \vspace{0.5em}
  \begin{enumerate}
    \item Abram \texttt{https://www.4devs.com.br/gerador\_de\_pessoas}
    \item Marquem para gerar o CPF \textbf{sem pontuação}
    \item Cliquem em \textbf{Gerar Pessoa}
    \item Anotem o nome e o CPF
    \item Façam isso três vezes
  \end{enumerate}

  \vspace{0.5em}
  O CPF gerado é válido, mas não é de ninguém.\\
  Por isso não usamos o CPF de vocês.
\end{frame}

\begin{frame}{Roteiro do teste}
  \begin{enumerate}
    \item Cadastrem as três pessoas, com senhas simples
    \item Entrem com a primeira e depositem R\$ 1.000
    \item Mandem R\$ 300 de PIX para a segunda
    \item Mandem R\$ 200 de PIX para a terceira
    \item Saquem R\$ 50
    \item Saiam e entrem com a segunda: saldo R\$ 300
    \item Da segunda, mandem R\$ 100 para a terceira
    \item Entrem com a terceira: saldo R\$ 300, dois PIX recebidos
  \end{enumerate}

\end{frame}

\begin{frame}[fragile]{Conferindo pelo phpMyAdmin}
\begin{lstlisting}[language=SQL,basicstyle=\ttfamily\footnotesize]
SELECT nome, saldo FROM conta;

SELECT * FROM movimentacao;

SELECT SUM(saldo) FROM conta;
\end{lstlisting}

  \vspace{0.4em}
  O último soma o saldo de todas as contas.\\
  Depois do roteiro, deve dar R\$ 950: os 1.000 depositados\\
  menos os 50 sacados.

  \vspace{0.4em}
  PIX não cria nem some dinheiro. Só muda de lugar.
\end{frame}

\begin{frame}{Critérios de aceite}
  \begin{itemize}
    \item [ ] Cadastro grava no MySQL e recusa CPF repetido
    \item [ ] Login recusa senha errada
    \item [ ] \texttt{/conta} sem login volta para o login
    \item [ ] Saldo aparece como \texttt{R\$ 1.234,56}
    \item [ ] Saque e PIX recusam quando falta saldo
    \item [ ] PIX recusa CPF que não existe e o próprio CPF
    \item [ ] Extrato mostra tudo, o mais novo em cima, com cores
    \item [ ] Reiniciar o Spring não perde nada
  \end{itemize}
\end{frame}

\begin{frame}{Erros que mais aparecem}
  \begin{itemize}
    \item Tabela \texttt{movimentacao} sem a coluna \texttt{conta\_id} ---
          faltou o \texttt{@ManyToOne}
    \item Extrato vazio sempre --- faltou chamar \texttt{registrar}
          ou faltou o \texttt{addAttribute("movimentacoes", ...)}
    \item \texttt{Exception evaluating SpringEL: isEmpty} --- o GET
          \texttt{/conta} não manda \texttt{movimentacoes}
    \item Valores do PIX invertidos --- o menos ficou no
          \texttt{registrar} errado
    \item Data com formato estranho --- conferir o
          \texttt{\#temporals.format}
  \end{itemize}
\end{frame}

% #############################################################################
\section{Recapitulando}

\begin{frame}{O que vimos nesses dois dias}
  \begin{enumerate}
    \item Thymeleaf: o servidor devolve a página já pronta
    \item \texttt{th:text}, \texttt{th:if}, \texttt{th:each}, \texttt{th:action}
    \item Formulário comum com \texttt{@RequestParam} e redirect depois do POST
    \item MySQL no XAMPP e o phpMyAdmin para olhar dentro
    \item \texttt{@Entity}, Repository e \texttt{save} no lugar do \texttt{ArrayList}
    \item Sessão como crachá do login
    \item \texttt{@Transactional} para o PIX ser tudo ou nada
    \item \texttt{@ManyToOne} ligando movimentação e conta
  \end{enumerate}
\end{frame}

\begin{frame}{Do ATM ao banco}
  \begin{center}
  \begin{tabular}{ll}
    \toprule
    ATM (aula 07) & Banco (aula 08) \\
    \midrule
    uma conta só & várias contas, com login \\
    saldo numa variável & saldo no MySQL \\
    parou o Spring, zerou & parou o Spring, continua \\
    JavaScript monta a tela & Thymeleaf monta a tela \\
    Axios chama \texttt{/api/conta} & o formulário chama \texttt{/deposito} \\
    sem histórico & extrato completo \\
    \bottomrule
  \end{tabular}
  \end{center}
\end{frame}

\begin{frame}{Se sobrar tempo}
  \begin{enumerate}
    \item Mostrar o CPF da pessoa embaixo do nome, na página da conta
    \item Não deixar sacar mais de R\$ 1.000 de uma vez
    \item Na tela de cadastro, recusar CPF que não tenha 11 caracteres
  \end{enumerate}

  \vspace{0.6em}
  Os três usam só o que vimos hoje:\\
  um \texttt{th:text}, um \texttt{if} no Controller, um \texttt{length()}.
\end{frame}

\begin{frame}
  \begin{center}
    \Huge Obrigado!\\[1.2em]
    \normalsize Mauricio Duailibi Neto\\
    Turma Java Entra21
  \end{center}
\end{frame}

\end{document}
